\section{Remaining Work}
\label{sec:future}

To finish the capstone vessel, the team must complete the following:
\begin{itemize}
\item Merge the \texttt{ltp} branch into \texttt{master} and sail the test
route on the water. The field test showed that the heading PID alone stalls when a
turn passes through the wind.
\item Enable the link-loss timeout and confirm the safe tail angle before
unattended operation.
\item Bring the radio link into compliance with FCC Part~15, by adding
frequency hopping or reducing output power (Section~\ref{sec:impact}).
\item Report each command's result individually, for example as a short
history in telemetry. In the field test, 27 of 65 command results were
overwritten before the ground station could see them.
\item Find the cause of the restart at 16:04:18 in the field test.
\item Confirm on the water that the \texttt{ltp} firmware's longer battery
window (20 readings instead of 2) removes the cell-voltage noise seen in the
field test, where individual cells varied by about \SI{180}{\milli\volt}
between packets.
\item Calibrate the simulator's hull drag and tail-to-angle-of-attack gain
from the field-test log, then retune both PIDs and the sail-planner table.
\item Run the full waypoint mission and record the data for
Section~\ref{sec:test}.
\end{itemize}
Before any longer deployment, the following changes are recommended:
\begin{itemize}
\item Replace the solderless breadboard with soldered connections.
\item Replace zinc-plated fasteners on the keel with stainless hardware.
\item Test the monohull's stability in rough water, and add outriggers if the
keel alone does not keep heel within limits.
\item Emulate the CM4 and radio in the simulator so the full command path can
be tested off the water.
\item Add an energy source, such as a solar panel, and use the cell-voltage
telemetry to manage energy on long missions.
\item Add obstacle avoidance and weather routing~\cite{stelzer2012}.
\end{itemize}